% Generated from locale/tr/content/first-order-logic/tableaux/proof-theoretic-notions.tex; do not hand-edit.


\iftag{FOL}
      {\olfileid[tr]{fol}{tab}{ptn}}
      {\olfileid[tr]{pl}{tab}{ptn}}
      
\olsection{İspat Kuramsal Kavramlar}

\begin{editorial}
Bu bölüm, tableau ağaçları için ispatlanabilirlik bağıntısı ile tutarlılık
tanımlarını bir araya getirir.
\end{editorial}

\begin{explain}
Bir dizi önemli semantik kavramı (geçerlilik, mantıksal gerektirme,
sağlanabilirlik) tanımladığımız gibi, şimdi bunlara karşılık gelen
\emph{ispat kuramsal kavramları} tanımlıyoruz. Bunlar, !!{sentence}s
ifadelerinin \iftag{FOL}{!!{structure}s}{değerlemeler} içinde sağlanmasına başvurularak değil, belirli
kapalı !!{tableau}s varlığına başvurularak tanımlanır. Bu kavramların
örtüştüğünün keşfedilmesi önemliydi. Örtüştüklerini söyleyen sonuçlar
\emph{sağlamlık} ve \emph{tamlık teoremleridir}.
\end{explain}


\begin{defn}[Teoremler]
Bir !!{sentence}~$!A$, $\sFmla{\False}{!A}$ için kapalı bir !!{tableau}
varsa \emph{teorem}dir. $!A$ teoremse $\Proves !A$, değilse $\Proves/ !A$
yazarız.
\end{defn}

\begin{defn}[!!^{derivability}]
Bir !!^a{sentence} için \emph{!!{derivable}} terimini şöyle kullanırız:
$!A$, bir !!{sentence}s kümesi~$\Gamma$'dan ancak ve ancak sonlu bir
$\{!B_1, \dots, !B_n\} \subseteq \Gamma$ kümesi ve
\[
\{
  \sFmla{\False}{!A},
  \sFmla{\True}{!B_1}, \dots,
  \sFmla{\True}{!B_n}
  \}.
\]
kümesi için kapalı bir !!{tableau} varsa türetilebilirdir; bunu
$\Gamma \Proves !A$ ile gösteririz. $!A$, $\Gamma$'dan !!{derivable}
değilse $\Gamma \Proves/ !A$ yazarız.
\end{defn}

\begin{defn}[Tutarlılık]
Bir !!{sentence}s kümesi~$\Gamma$, ancak ve ancak sonlu bir
$\{!B_1, \dots, !B_n\} \subseteq \Gamma$ kümesi ve
\[
\{
  \sFmla{\True}{!B_1}, \dots,
  \sFmla{\True}{!B_n}  
  \}.
\]
kümesi için kapalı bir !!{tableau} varsa \emph{tutarsızdır}. $\Gamma$
tutarsız değilse ona \emph{tutarlı} deriz.
\end{defn}

\begin{prop}[Yansımalılık]
\ollabel{prop:reflexivity}
$!A \in \Gamma$ ise $\Gamma \Proves !A$.
\end{prop}

\begin{proof}
$!A \in \Gamma$ ise $\{!A\}$,~$\Gamma$'nın sonlu bir alt kümesidir ve şu
!!{tableau} kapalıdır:
\begin{oltableau}
  [\sFmla{\False}{\formula{A}}, just = \TAss
    [\sFmla{\True}{\formula{A}}, just = \TAss,close]
  ]
\end{oltableau}
\end{proof}
  

\begin{prop}[Monotonluk]
\ollabel{prop:monotonicity}
$\Gamma \subseteq \Delta$ ve $\Gamma \Proves !A$ ise $\Delta
\Proves !A$.
\end{prop}

\begin{proof}
~$\Gamma$'nın her sonlu alt kümesi aynı zamanda~$\Delta$'nın sonlu bir alt
kümesidir.
\end{proof}

\begin{prop}[Geçişlilik]
\ollabel{prop:transitivity}
$\Gamma \Proves !A$ ve $\{!A\} \cup
\Delta \Proves !B$ ise $\Gamma \cup \Delta \Proves !B$.
\end{prop}

\begin{proof}
$\{!A\} \cup \Delta \Proves !B$ ise, aşağıdaki kümenin kapalı bir
!!{tableau}'su olacak biçimde sonlu bir $\Delta_0 =
\{!C_1, \dots, !C_n\} \subseteq \Delta$ alt kümesi vardır:
\begin{align*}
\{\sFmla{\False}{!B}, & \sFmla{\True}{!A}, \sFmla{\True}{!C_1},
\dots, \sFmla{\True}{!C_n}\}.
\intertext{$\Gamma \Proves !A$ ise, aşağıdaki kümenin kapalı bir
  !!{tableau}'su olacak biçimde $\{!D_1, \dots, !D_m\} \subseteq \Gamma$
  vardır:}
\{\sFmla{\False}{!A}, & \sFmla{\True}{!D_1},
\dots, \sFmla{\True}{!D_m}\}.
\end{align*}

Şimdi varsayımları
\[
\sFmla{\False}{!B},
\sFmla{\True}{!C_1}, \dots, \sFmla{\True}{!C_n},
\sFmla{\True}{!D_1}, \dots, \sFmla{\True}{!D_m}.
\]
olan !!{tableau}'yu ele alın. $!A$ üzerinde \Cut{} kuralını uygulayın. Bu,
biri $\sFmla{\True}{!A}$, diğeri $\sFmla{\False}{!A}$ içeren iki dal
oluşturur. Dolayısıyla dallardan birinde
\[
\{\sFmla{\False}{!B}, \sFmla{\True}{!A},
\sFmla{\True}{!C_1}, \dots, \sFmla{\True}{!C_n}\}
\]
ifadelerinin tümü bulunur. Bu varsayımlar için kapalı bir !!{tableau}
olduğundan onu bu dala ekleyebiliriz; $\sFmla{\True}{!A}$ içinden geçen her
dal kapanır. Diğer dalda
\[
\{\sFmla{\False}{!A}, \sFmla{\True}{!D_1}, \dots,
\sFmla{\True}{!D_m}\}
\]
ifadelerinin tümü bulunur; bu nedenle kapalı bir !!{tableau} elde etmek için
öteki tarafı da tamamlayabiliriz. Böylece $\Gamma \cup \Delta \Proves !B$
olduğu gösterilir.
\end{proof}

Bunun özellikle şu anlama geldiğine dikkat edin: $\Gamma \Proves !A$ ve
$!A \Proves !B$ ise $\Gamma \Proves !B$. Ayrıca $!A_1,
\dots, !A_n \Proves !B$ ve her~$i$ için $\Gamma \Proves !A_i$ ise
$\Gamma \Proves !B$ sonucu çıkar.

\begin{prop}
\ollabel{prop:incons}
$\Gamma$ ancak ve ancak her !!{sentence}~$!A$ için $\Gamma \Proves !A$
ise tutarsızdır.
\end{prop}

\begin{proof}
Alıştırma.
\end{proof}

\tagprob{FOL}
\begin{prob}
\olref[fol][tab][ptn]{prop:incons} önermesini ispatlayın.
\end{prob}
\tagendprob

\tagprob{notFOL}
\begin{prob}
\olref[pl][tab][ptn]{prop:incons} önermesini ispatlayın.
\end{prob}
\tagendprob

\begin{prop}[Kompaktlık]
  \ollabel{prop:proves-compact}
  \begin{enumerate}
  \item $\Gamma \Proves !A$ ise $\Gamma_0 \Proves !A$ olacak biçimde sonlu
    bir $\Gamma_0 \subseteq \Gamma$ alt kümesi vardır.
  \item ~$\Gamma$'nın her sonlu alt kümesi tutarlıysa $\Gamma$ tutarlıdır.
  \end{enumerate}
\end{prop}

\begin{proof}
  \begin{enumerate}
    \item $\Gamma \Proves !A$ ise sonlu bir
      $\Gamma_0 = \{!B_1, \dots, !B_n\} \subseteq \Gamma$ alt kümesi ve
      \[
      \{\sFmla{\False}{!A}, \sFmla{\True}{!B_1}, \dots, \sFmla{\True}{!B_n}\}
      \]
      için kapalı bir !!{tableau} vardır. Bu !!{tableau} aynı zamanda
      $\Gamma_0 \Proves !A$ olduğunu gösterir.
    \item $\Gamma$ tutarsızsa sonlu bir
      $\Gamma_0 = \{!B_1, \dots, !B_n\}$ alt kümesi için
      \[
      \{\sFmla{\True}{!B_1}, \dots, \sFmla{\True}{!B_n}\}
      \]
      kümesinin kapalı bir !!{tableau}'su vardır. Bu kapalı !!{tableau},
      $\Gamma_0$'ın tutarsız olduğunu gösterir.
  \end{enumerate}
\end{proof}

